
\subsection{$\mathcal M_{k+1}^{\mathrm{main,reg}}(\beta)$ 上的仓西坐标图：弗雷德霍姆理论}
\label{subsec:7-1-2}

本小节及第~7.1.3 小节将给出命题~\ref{prop:7-1-1}
证明中的分析细节\liftIndex{I227}\liftIndex{I228}。
为此，需要在
$\mathcal M^{\mathrm{main,reg}}_{k+1}$
的元素 $\Sigma$ 上固定度量。
我们使用第~3.4 节引入的映射
$\Theta\colon\operatorname{Gr}_k\to
\mathcal M^{\mathrm{main,reg}}_k$；
在那里已为
$\mathcal M^{\mathrm{main,reg}}_k$
的各元素固定了一族光滑变化的黎曼度量。

现固定一个与 $\omega$ 相容的几乎复结构，
相应的柯西--黎曼算子记为 $\bar\partial$。
设 $L\subset(M,\omega)$ 为拉格朗日子流形。
取一对 $((\Sigma,\vec z),w)$，其中
$(\Sigma,\vec z)\in\mathcal M^{\mathrm{main,reg}}_k$，
$w\colon(\Sigma,\partial\Sigma)\to(M,L)$ 为光滑映射。
以 $\beta\in\Pi(M,L)$ 表示 $w$ 的同伦类的等价类。
如第~3.4 节，将
\[
\Sigma\setminus\{z_1,\ldots,z_k\}\cong X(\mathfrak t,\ell),
\]
从而把 $w$ 视为如下映射：
\[
w\colon(X(\mathfrak t,\ell),\partial X(\mathfrak t,\ell))
\longrightarrow(M,L).
\]
反过来，给定这样的 $w$，
若要将其紧化为
$(\Sigma,\partial\Sigma)\to(M,L)$ 上的映射，
便须限制 $w$ 的渐近行为。
由此考虑下述巴拿赫流形。
取并固定充分小的 $\delta>0$，再取充分大的 $p>0$。

\begin{definition}{7.1.3}\label{def:7-1-3}
$W_k^{1,p}(M,L)$ 是满足下列条件的全部
$((\mathfrak t,\ell),w,\vec p)$ 的集合：
\begin{align*}
& (\mathfrak t,\ell)\in\operatorname{Gr}_k;
  \tag{7.1.4.1}\label{eq:7-1-4-1}\\
& w\colon(X(\mathfrak t,\ell),\partial X(\mathfrak t,\ell))
   \to(M,L)\text{ 属于 }W^{1,p}_{\mathrm{loc}},
   \text{ 即局部属于 }W^{1,p};
  \tag{7.1.4.2}\label{eq:7-1-4-2}\\
& \vec p=(p_1,\ldots,p_k),\quad p_i\in L.
  \tag{7.1.4.3}\label{eq:7-1-4-3}
\end{align*}
此外，取任意 $e\in C^1_{\mathrm{ext}}(\mathfrak t)$，
并将
$L_e\cong(-\infty,0]\times[0,1]
\subset X(\mathfrak t,\ell)$；
以 $\tau$ 表示 $(-\infty,0]$ 的参数，
以 $t$ 表示 $[0,1]$ 的参数，则要求
\begin{equation}
\begin{aligned}
&\int_{L_e}e^{\delta|\tau|'}|\nabla w|^p\,d\tau\,dt\\
&\qquad+\int_{L_e}e^{\delta|\tau|'}
\operatorname{dist}(w(\tau,t),p_i)^p\,d\tau\,dt<\infty.
\end{aligned}
\tag{7.1.4.4}\label{eq:7-1-4-4}
\end{equation}
其中 $\operatorname{dist}$ 是黎曼距离。
在这里及下文中，
$|\cdot|'\colon\mathbb R\to[0,\infty)$
表示一个光滑函数，满足当 $|\tau|\ge1$ 时
$|\tau|'=|\tau|$，当 $|\tau|\le1/2$ 时
$|\tau|'=0$。
\end{definition}

在下一引理及以后，
取 $TM$ 的一个幺正联络，使平行移动保持 $TL$，
并用它定义式~(7.1.6.5) 等处的同构
\[
\operatorname{Pal}\colon T_{p_i}M\longrightarrow T_{w(\tau,t)}M.
\]

\begin{lemma}{7.1.5}\label{lem:7-1-5}
空间 $W_k^{1,p}(M,L)$ 是巴拿赫流形。
在 $((\mathfrak t,\ell),w,\vec p)$ 处，
切空间
$T_{((\mathfrak t,\ell),w,\vec p)}W_k^{1,p}(M,L)$
可与满足下列条件的 $(\xi,V,\vec v)$ 的集合等同：
\begin{align*}
&\xi\in T_{(\mathfrak t,\ell)}\operatorname{Gr}_k;
  \tag{7.1.6.1}\label{eq:7-1-6-1}\\
&V\text{ 是 }w^*(TM)\text{ 的 }W^{1,p}_{\mathrm{loc}}\text{ 截面};
  \tag{7.1.6.2}\label{eq:7-1-6-2}\\
&\vec v=(v_1,\ldots,v_k),\quad v_i\in T_{p_i}L;
  \tag{7.1.6.3}\label{eq:7-1-6-3}\\
&V(z)\in T_{w(z)}L\quad
  (z\in\partial X(\mathfrak t,\ell)).
  \tag{7.1.6.4}\label{eq:7-1-6-4}
\end{align*}
若 $e\in C^1_{\mathrm{ext}}(\mathfrak t)$，
$L_e\cong(-\infty,0]\times[0,1]
\subset X(\mathfrak t,\ell)$，还须有
\begin{equation}
\int_{L_e}e^{\delta|\tau|'}
\bigl(|\nabla(V-\operatorname{Pal}v_i)|^p
      +|V-\operatorname{Pal}v_i|^p\bigr)
\,d\tau\,dt<\infty.
\tag{7.1.6.5}\label{eq:7-1-6-5}
\end{equation}
这里
$\operatorname{Pal}\colon T_{p_i}M\to T_{w(\tau,t)}M$
是平行移动。
$W_k^{1,p}(M,L)$ 上的范数定义为
\[
\begin{aligned}
\|(\xi,V,\vec v)\|_{1,p}^p
={}&\|\xi\|^p+|\vec v|^p
+\sum_{e\in C^1_{\mathrm{int}}(\mathfrak t)}
  \int_{L_e}(|\nabla V|^p+|V|^p)\,d\tau\,dt\\
&+\sum_{e_i\in C^1_{\mathrm{ext}}(\mathfrak t)}
  \int_{L_{e_i}}e^{\delta|\tau|'}
  \bigl(|\nabla(V-\operatorname{Pal}v_i)|^p
       +|V-\operatorname{Pal}v_i|^p\bigr)
  \,d\tau\,dt.
\end{aligned}
\]
\end{lemma}

证明是标准的，这里从略。
注意上面的 $\vec v$ 由 $V$ 唯一决定。

接下来在 $W_k^{1,p}(M,L)$ 上定义巴拿赫丛
$\mathcal E^{0,p}$。
取 $((\mathfrak t,\ell),w,\vec p)
\in W_k^{0,p}(M,L)$，
考虑 $\Lambda^{0,1}\otimes w^*(TM)$ 的截面 $V$，
满足：
\begin{align*}
&V\in W^{0,p}_{\mathrm{loc}};
  \tag{7.1.7.1}\label{eq:7-1-7-1}\\
&\text{对 }e\in C^1_{\mathrm{ext}}(\mathfrak t),\quad
  L_e\cong(-\infty,0]\times[0,1]
  \subset X(\mathfrak t,\ell),\quad
  \int_{L_e}e^{\delta|\tau|'}|V|^p
  \,d\tau\,dt<\infty.
  \tag{7.1.7.2}\label{eq:7-1-7-2}
\end{align*}
把满足 (7.1.7.1)--(7.1.7.2) 的全部 $V$
记为 $\mathcal E^{0,p}_{((\mathfrak t,\ell),w,\vec p)}$。
这是一个巴拿赫空间，其范数由下式定义：
\begin{equation}
\begin{aligned}
\|V\|_{0,p}^p
={}&\sum_{e\in C^1_{\mathrm{int}}(\mathfrak t)}
\int_{L_e}|V|^p\,d\tau\,dt\\
&+\sum_{e_i\in C^1_{\mathrm{ext}}(\mathfrak t)}
\int_{L_{e_i}}e^{\delta|\tau|'}|V|^p\,d\tau\,dt.
\end{aligned}
\tag{7.1.7.3}\label{eq:7-1-7-3}
\end{equation}
给并集
\[
\mathcal E^{0,p}:=
\bigcup_{((\mathfrak t,\ell),w,\vec p)}
\mathcal E^{0,p}_{((\mathfrak t,\ell),w,\vec p)}
\]
赋予显然的拓扑后，
$\mathcal E^{0,p}$ 成为 $W_k^{1,p}(M,L)$
上的巴拿赫丛。
事实上，此丛可延拓到 $W_k^{0,p}(M,L)$ 上。
容易证明下列结论。

\begin{lemma}{7.1.8}\label{lem:7-1-8}
对每个
$((\mathfrak t,\ell),w,\vec p)\in W_k^{1,p}(M,L)$，
都有
$\bar\partial w\in
\mathcal E^{0,p}_{((\mathfrak t,\ell),w,\vec p)}$。
此外，映射
$((\mathfrak t,\ell),w,\vec p)\mapsto\bar\partial w$
定义了 $\mathcal E^{0,p}$ 的一个光滑截面。
\end{lemma}

利用沿最短测地线的平行移动，
可给 $\mathcal E^{0,p}$ 定义局部平凡化。
这里选择度量，使 $L$ 是全测地的。
此处及下文均取 $TM$ 的幺正联络，
使平行移动保持 $TL$，
并用它定义上述同构。
在这一平凡化下，定义 $\bar\partial$ 的协变导数
\begin{equation}
D_{((\mathfrak t,\ell),w,\vec p)}\bar\partial:
T_{((\mathfrak t,\ell),w,\vec p)}W_k^{1,p}(M,L)
\longrightarrow
\mathcal E^{0,p}_{((\mathfrak t,\ell),w,\vec p)}.
\tag{7.1.9}\label{eq:7-1-9}
\end{equation}

\begin{lemma}{7.1.10}\label{lem:7-1-10}
若 $\delta>0$ 充分小，
则 $D_{((\mathfrak t,\ell),w,\vec p)}\bar\partial$
是弗雷德霍姆算子，其弗雷德霍姆指标为
$\mu_L(\beta)+k-3+n$。
这里 $\beta\in\pi_2(M,L)$ 是 $w$ 的同伦类，
$\mu_L(\beta)$ 是它的马斯洛夫指标。
\end{lemma}

证明现在已属标准论证，故略。
注意这里需要权重 $e^{\delta|\tau|'}$
使 $D_{((\mathfrak t,\ell),w,\vec p)}\bar\partial$
成为弗雷德霍姆算子，
因为该算子在端部退化。

下面的结论是指数衰减估计的推论
（见后述引理~7.1.37）。

\begin{lemma}{7.1.11}\label{lem:7-1-11}
设 $((\Sigma,\vec z),w)$ 满足
$(\Sigma,\vec z)\in\mathcal M_k^{\mathrm{main,reg}}$，
且 $w\colon(\Sigma,\partial\Sigma)\to(M,L)$
局部属于 $L^{1,p}$ 类。
假设 $\bar\partial w$ 具有紧支集。
将 $\Sigma\setminus\{z_1,\ldots,z_k\}$
与 $X(\mathfrak t,\ell)$ 等同，令 $p_i=w(z_i)$。
若 $\delta$ 充分小，则
$((\mathfrak t,\ell),w,\vec p)\in W_k^{1,p}(M,L)$。
\end{lemma}

现取
$((\mathfrak t,\ell),w,\vec p)\in W_k^{1,p}(M,L)$
且 $\bar\partial w=0$。
原书称，由引理 29.11 可知，
此元素对应于某个 $\beta$ 的
$\mathcal M_k^{\mathrm{main}}(\beta)$
中的元素。
\translatornote{源书实体页 412 引用“引理 29.11”，
但锁定本书没有该引理。
本节引理 7.1.11 的方向是从已定义在紧化曲面上的映射
推出加权条件，不能单独证明反向的可去奇性与延拓。
这里保留源书引据及结论，不猜改编号；见疑误 A-0408。}

\begin{proposition}{7.1.12}\label{prop:7-1-12}
在空间 $\mathcal M_k^{\mathrm{main}}(\beta)$ 中，
$((\mathfrak t,\ell),w,\vec p)$
存在一个仓西邻域。
\end{proposition}

\begin{proof}
考虑算子~(7.1.9)。
由引理~\ref{lem:7-1-10}，
可取
$\mathcal E^{0,p}_{((\mathfrak t,\ell),w,\vec p)}$
的一个有限维子空间
$E_{((\mathfrak t,\ell),w,\vec p)}$，
使下列映射满射：
\begin{equation}
\begin{aligned}
\pi\circ D_{((\mathfrak t,\ell),w,\vec p)}\bar\partial:
&\ T_{((\mathfrak t,\ell),w,\vec p)}W_k^{1,p}(M,L)\\
&\longrightarrow
\mathcal E^{0,p}_{((\mathfrak t,\ell),w,\vec p)}
/E_{((\mathfrak t,\ell),w,\vec p)}.
\end{aligned}
\tag{7.1.13}\label{eq:7-1-13}
\end{equation}
其中 $\pi$ 是自然投影。
由唯一延拓定理 \textbf{[Aro57]}，
还可选择 $E_{((\mathfrak t,\ell),w,\vec p)}$，
使其中每个 $V$ 满足：
\begin{align*}
&V\text{ 光滑};
  \tag{7.1.14.1}\label{eq:7-1-14-1}\\
&V\text{ 具有紧支集}.
  \tag{7.1.14.2}\label{eq:7-1-14-2}
\end{align*}

取 $(\mathfrak t,\ell)$ 在 $\operatorname{Gr}_k$
中的小邻域 $\mathcal U$。
对 $(\mathfrak t',\ell')\in\mathcal U$，
可取如下微分同胚：
\[
I_{(\mathfrak t,\ell),(\mathfrak t',\ell')}:
X(\mathfrak t,\ell)\longrightarrow X(\mathfrak t',\ell')
\]
它具有下列性质：
\begin{align*}
&I_{(\mathfrak t,\ell),(\mathfrak t',\ell')}
\text{ 光滑地依赖于 }(\mathfrak t',\ell');
\tag{7.1.15.1}\label{eq:7-1-15-1}\\
&I_{(\mathfrak t,\ell),(\mathfrak t',\ell')}
\text{ 在某个紧集之外为恒等映射}.
\tag{7.1.15.2}\label{eq:7-1-15-2}
\end{align*}
在 (7.1.15.1) 中，
相对于 $(\mathfrak t',\ell')$ 的光滑性
由 $\mathcal U$ 上开黎曼曲面的万有族来定义。
在 (7.1.15.2) 中，
$X(\mathfrak t,\ell)$ 去掉一个紧集后，
可与各 $e\in C^1_{\mathrm{ext}}(\mathfrak t)$
所对应的
$(-\infty,-T]\times[0,1]\subset L_e$
之并等同。

现在取
$((\mathfrak t',\ell'),w',\vec p')\in W_k^{1,p}(M,L)$，
它位于 $((\mathfrak t,\ell),w,\vec p)$ 的一个小邻域中。
沿最短测地线平行移动，
得到典范同构
\[
\mathcal E^{0,p}_{((\mathfrak t,\ell),w,\vec p)}
\cong
\mathcal E^{0,p}_{((\mathfrak t',\ell'),w',\vec p')}.
\]
这里仍取平行移动保持 $TL$ 的 $TM$ 幺正联络
来定义此同构；下文不再重复说明。
考虑方程
\begin{equation}
\bar\partial w'\equiv0
\pmod{E_{((\mathfrak t,\ell),w,\vec p)}}.
\tag{7.1.16}\label{eq:7-1-16}
\end{equation}
这里通过上述同构，
将右端的 $E_{((\mathfrak t,\ell),w,\vec p)}$
视为
$\mathcal E^{0,p}_{((\mathfrak t',\ell'),w',\vec p')}$
的子空间。
\translatornote{源书实体页 412 在这里写成基点处的
$\mathcal E^{0,p}_{((\mathfrak t,\ell),w,\vec p)}$；
但右端 $E_{((\mathfrak t,\ell),w,\vec p)}$
本已是该纤维的子空间。
式 (7.1.16) 用于邻点 $w'$，故须经上述平行移动
把它视为邻点纤维的子空间；见疑误 A-0409。}

\begin{lemma}{7.1.17}\label{lem:7-1-17}
令 $U((\mathfrak t,\ell),w,\vec p)$
为满足 (7.1.16) 的全部
$((\mathfrak t',\ell'),w',\vec p')$ 的集合。
则它是一个光滑流形，维数为
\[
\mu_L(\beta)+k-3+n
+\dim E_{((\mathfrak t,\ell),w,\vec p)}.
\]
\end{lemma}

这一结论直接由引理~\ref{lem:7-1-10}
和隐函数定理得出。
定义如下映射：
\[
\begin{gathered}
s_{((\mathfrak t,\ell),w,\vec p)}:
U((\mathfrak t,\ell),w,\vec p)
\longrightarrow E_{((\mathfrak t,\ell),w,\vec p)},\\
s_{((\mathfrak t,\ell),w,\vec p)}
((\mathfrak t',\ell'),w',\vec p')=\bar\partial w'.
\end{gathered}
\]
由 $U((\mathfrak t,\ell),w,\vec p)$ 的定义，
右端属于 $E_{((\mathfrak t,\ell),w,\vec p)}$。
回忆当 $k>0$ 时，
$(\Sigma,\vec z)\in
\mathcal M_k^{\mathrm{main,reg}}(\beta)$
的自同构群无挠。
因此 $((\mathfrak t,\ell),w,\vec p)$
的自同构群平凡，可直接取
$\Gamma_{((\mathfrak t,\ell),w,\vec p)}=\{1\}$。
只有当内部出现球面气泡时，
才需要非平凡的 $\Gamma$。
这完成了命题~\ref{prop:7-1-12} 的证明。
\end{proof}

不过，关于评价映射还须对仓西邻域施加一个条件。
评价映射定义如下：
\[
\begin{gathered}
\operatorname{ev}:
U((\mathfrak t,\ell),w,\vec p)
\longrightarrow L^k,\\
\operatorname{ev}
((\mathfrak t',\ell'),w',(p'_1,\ldots,p'_k))
=(p'_1,\ldots,p'_k).
\end{gathered}
\]

\begin{lemma}{7.1.18}\label{lem:7-1-18}
可选择命题~\ref{prop:7-1-12} 中的仓西邻域，
使 $\operatorname{ev}$ 成为浸没。
\end{lemma}

\begin{proof}
考虑
$T_{((\mathfrak t,\ell),w,\vec p)}W_k^{1,p}(M,L)$
中由 $\vec v=0$ 的
$(\xi,V,\vec v)$ 组成的子空间，
记作 $(d\operatorname{ev})^{-1}(0)$。
将算子~(7.1.9) 限制于此，得到
\[
D_{((\mathfrak t,\ell),w,\vec p)}\bar\partial
\big|_{(d\operatorname{ev})^{-1}(0)}:
(d\operatorname{ev})^{-1}(0)
\longrightarrow
\mathcal E^{0,p}_{((\mathfrak t,\ell),w,\vec p)}.
\]
注意 $(d\operatorname{ev})^{-1}(0)$
在
$T_{((\mathfrak t,\ell),w,\vec p)}W_k^{1,p}(M,L)$
中具有有限余维，
故这一限制仍是弗雷德霍姆算子。
于是可选择
$E_{((\mathfrak t,\ell),w,\vec p)}$，
使
\[
\begin{aligned}
\pi\circ D_{((\mathfrak t,\ell),w,\vec p)}\bar\partial:
&(d\operatorname{ev})^{-1}(0)\\
&\longrightarrow
\mathcal E^{0,p}_{((\mathfrak t,\ell),w,\vec p)}
/E_{((\mathfrak t,\ell),w,\vec p)}
\end{aligned}
\]
满射。
由此容易看出，
对这个 $E_{((\mathfrak t,\ell),w,\vec p)}$，
引理~\ref{lem:7-1-18} 的结论成立。
\end{proof}
